\section{Hướng dẫn sử dụng}

\subsection{Đăng ký, đăng nhập và đăng xuất}
Mở trang đăng ký, nhập email, mật khẩu dài 8--128 ký tự và xác nhận trùng mật khẩu,
sau đó gửi form. Đăng ký thành công sẽ chuyển sang đăng nhập.
Nhập email và mật khẩu để mở không gian công việc.
Muốn đăng xuất, mở menu Tài khoản ở góc trên bên phải và chọn Đăng xuất.
Phiên hết hạn hoặc backend khởi động lại sẽ yêu cầu đăng nhập lại.
\begin{itemize}
  \item \#TODO: Thêm ảnh chụp màn hình minh họa.
\end{itemize}

\subsection{Công việc và dự án}
\begin{enumerate}
  \item Chọn Công việc. Nhập tên ở ô thêm nhanh để tạo việc với bộ lọc hiện tại;
  chọn thêm công việc để nhập đầy đủ mô tả, dự án, hạn, ưu tiên và trạng thái.
  Tên việc bắt buộc, tối đa 200 ký tự; mô tả tối đa 2000 ký tự.
  \item Chọn Không gắn ưu tiên nếu không cần tag; chọn Cao hoặc Thấp để đánh dấu.
  Muốn bỏ tag của việc cũ, mở form sửa, chọn Không gắn ưu tiên rồi lưu.
  \item Bấm tên việc để sửa. Bấm nút tròn để hoàn thành/mở lại;
  dùng ô trạng thái để chuyển sang Đang làm. Menu công việc có thao tác xóa và yêu cầu xác nhận.
  \item Dùng thanh bên và bộ lọc để chọn dự án, hạn hoặc trạng thái;
  tìm kiếm theo tên/mô tả. Có 3 cách sắp xếp: Mới tạo, Hạn gần nhất, Hạn xa nhất.
  \item Thêm dự án, nhập tên dài 1--80 ký tự và chọn màu.
  Chọn dự án để xem công việc bên trong; menu dự án cho phép sửa/xóa.
  Xóa dự án giữ công việc và chuyển chúng về Chưa chia dự án.
\end{enumerate}

\subsection{Lịch biểu}
Chọn Lịch biểu, chuyển tháng và chọn ngày để xem công việc đến hạn.
Bấm công việc để sửa; chọn Thêm việc cho ngày này để đặt sẵn hạn theo ngày đang chọn.
Có thể lọc theo trạng thái. Công việc không có hạn không xuất hiện trên lịch.

\subsection{Ghi chú}
Chọn Ghi chú và tạo giấy mới. Nhập tiêu đề hoặc nội dung, chọn ngày và một trong 4 màu,
sau đó lưu. Tiêu đề tối đa 120 ký tự, nội dung tối đa 10000 ký tự.
Bấm giấy để sửa; kéo thanh đầu giấy vào khung trống để xếp ngay ngắn.
Vị trí trong khung tự lưu, giữ sau khi tải lại trang; ô có giấy không nhận chồng.
Kéo ra ngoài các khung để tiếp tục xem trên lớp nổi khi chuyển mục;
khung cũ vẫn để trống. Kéo vào khung hoặc nhấn X để cất giấy về bảng.
Giấy nổi là trạng thái của phiên giao diện;
việc lưu nội dung, ngày và màu được thực hiện qua nút lưu.
Xóa giấy cần xác nhận; đóng form có thay đổi chưa lưu sẽ hỏi trước khi bỏ.

\subsection{Tổng quan và ghi nhận kết quả}
Chọn Tổng quan để xem số việc hoàn thành trong 7 ngày (hôm nay và 6 ngày trước),
số đang làm, quá hạn, tỷ lệ hoàn thành toàn bộ công việc và 5 việc hoàn thành gần đây.
Ngày thống kê theo giờ Việt Nam. Chưa có công việc thì tỷ lệ hiển thị dấu gạch.
Sau khi hoàn thành việc, chọn Ghi nhận kết quả để viết ghi chú và tự đánh giá
Tốt, Đạt yêu cầu hoặc Cần cải thiện. Có thể bỏ qua bước này.
Mở lại việc sẽ xóa mốc hoàn thành và giữ ghi chú kết quả đã lưu.
\begin{itemize}
  \item \#TODO: Thêm ảnh chụp màn hình cho các thao tác quan trọng.
\end{itemize}
